% sec-00: abstract, executive summary, reader's guide, and the two indexes.
% FULL stage.  Every drafting instruction resolved.
\section*{Abstract}
\phantomsection\addcontentsline{toc}{section}{Abstract}

In August 2026 the author filed ten funding applications for a Phase 1
first-in-human trial of an LLM-advised robotic pancreaticoduodenectomy with
perioperative daraxonrasib, written throughout as an independent scientist
\cite{tenapps2026}. Nine of the ten address mechanisms that fund a person or an
institution. One, application 05, addresses the only mechanism in the set built
for a company, and it was the shortest of the ten. This paper converts the
author from independent scientist to small-business operator and rewrites that
application as the document it should have been.

The clause the conversion turns on is neither the report's individual-scientist
language nor its novel-performers chapter. It is the Small Business Innovation
Research clause of \emph{Science: A New Golden Age}, which appears in three
separate chapters and is the only clause in the report written for a firm
\cite{kratsios2026sciencegoldenage}. Three numbers govern what follows:
\$306,000 for a nine-month Phase I, \$1,300,000 for a twenty-four-month Phase
II, and the \$3,500,000 five-year direct programme the two are measured against
\cite{fundapp2}. Two more are derived: \$1,396,000 of direct work sits inside
the \$1,606,000 award, and \$2,104,000 sits outside it.

What is offered is a costed twelve-milestone schedule. Every row carries an
evidence artifact that exists as a file, a cost that sums into its phase, a stop
condition stated before the work begins, and the funder mechanism that pays.

\section*{Executive Summary}
\phantomsection\addcontentsline{toc}{section}{Executive Summary}

\textbf{The clause.} \emph{Science: A New Golden Age} names SBIR three times: in
the Chapter I summary of recommendations, in Chapter III on deploying SBIR and
STTR strategically to couple federally-seeded companies with the scientific
enterprise, and in Chapter IV on SBIR opening doors for technician-founded
ventures. Its novel-performers chapter describes work at tens of millions of
dollars with teams of ten to a hundred people over half a decade. This company
proposes \$1,606,000, 2.6 full-time equivalents, and 33 months.

\textbf{The entity.} ChemicalQDevice holds thirteen deposited assets outright
and none under encumbrance: two protocols, an IND package, a simulation stack, a
verification suite, and a public repository. It licenses three things, none
exclusively. It has contracted for nothing. Seven instruments are absent, and
not one can be obtained by the company alone.

\textbf{The two prices.} The identical four-layer budget costs \$3,500,000 of
direct work over sixty months and \$1,396,000 of direct work over thirty-three.
The award that carries the second is \$1,606,000, a 15.0 percent overhead and
fee load against an SBIR allowance of 40 percent that the company does not
claim. Routed through a university at a 57 percent facilities and administrative
rate, the same direct work would cost a funder \$2,137,000.

\textbf{The bridge.} Federal money of \$1,606,000 sits below a firewall;
contributed drug, theatre and pathology sit unpriced in the middle; and
\$5,900,000 of private capital sits above, raised only after a milestone closes.
That is 3.67 to one on cash alone, against the annexed FY 2028 target of at
least three to one \cite{kratsiosvought2026fy2028priorities}.

\textbf{The milestones.} Twelve rows, five in Phase I summing to \$306,000 and
seven in Phase II summing to \$1,300,000. Each is closed by three concurrent
evidence branches that must join before an audit can compare them.

\textbf{What the conversion costs.} The individual-scientist framing that made
all ten prior applications distinctive is abandoned. Capitalization detail the
author might prefer private is published permanently under a DOI. Three of the
ten prior recipients will read a company plan as off-mechanism. Those four
concessions are stated in \S9 and are not softened anywhere.

\section*{How to Read This Document}
\phantomsection\addcontentsline{toc}{section}{How to Read This Document}

The document is written so that four different readers can each open it at a
different page and find the thing they came for within two pages.

\begin{apptable}
\begin{tabularx}{\textwidth}{>{\raggedright\arraybackslash}p{3.4cm} >{\raggedright\arraybackslash}p{2.1cm} >{\raggedright\arraybackslash}p{2.4cm} >{\raggedright\arraybackslash}X}
\headrow \thead{Reader} & \thead{Open at} & \thead{Then} & \thead{The question this reader is actually asking}\\
\midrule
SBIR program director & \S3, Table 8 & \S5, Tables 13 and 14 & Does the award buy a defensible gate, and does the budget sum? \\
ARPA-H programme manager & \S5, Figure 13 & \S1, Figure 1 & Is there a milestone I can hold someone to, and a stop condition? \\
Private investor & \S4, Table 12 & \S2, Table 6 & What do I own, when may I own it, and what is already built? \\
UC San Diego intake & \S8, Figure 19 & \S7, Figure 18 & What is being asked of this institution, and by when? \\
\end{tabularx}
\end{apptable}
\tabcap{Four readers and the section each should open first. No reader is asked
to read the document front to back, and every entry point is a table or a figure
rather than a paragraph.}

\begin{apptable}
\begin{tabularx}{\textwidth}{>{\raggedright\arraybackslash}p{0.7cm} >{\raggedright\arraybackslash}p{1.1cm} >{\raggedright\arraybackslash}p{2.5cm} >{\raggedright\arraybackslash}p{3.3cm} >{\raggedright\arraybackslash}X}
\headrow \thead{Fig} & \thead{\S} & \thead{Platform} & \thead{Source directory} & \thead{The question it answers}\\
\midrule
1 & 1 & Mermaid & \appfile{mermaid/} & Which clause does a company with no indirect-cost base qualify under? \\
2 & 1 & D2 & \appfile{d2/} & What kind of institution is a 2.6 FTE firm, on the report's own axes? \\
3 & 1 & Graphviz & \appfile{graphviz/} & What does the same direct work cost under three overhead regimes? \\
4 & 2 & Diagrams & \appfile{diagrams-python/} & What is the shape of what the company holds today? \\
5 & 2 & D2 & \appfile{d2/} & What is owned, licensed, contracted, and absent, with terms? \\
6 & 2 & PlantUML & \appfile{plantuml/} & What may the sponsor do alone, and what needs an institution? \\
7 & 3 & Mermaid & \appfile{mermaid/} & What must hold at month nine for Phase II to begin? \\
8 & 3 & D2 & \appfile{d2/} & What does the same programme cost under two mechanisms? \\
9 & 3 & Graphviz & \appfile{graphviz/} & Which work packages does the money reach, and which not? \\
10 & 4 & PlantUML & \appfile{plantuml/} & Which capital positions may be held while a participant is on study? \\
11 & 4 & D2 & \appfile{d2/} & Where does every dollar come from, and what separates the tiers? \\
12 & 4 & Mermaid & \appfile{mermaid/} & In what order is a private round executed during enrollment? \\
13 & 5 & Mermaid & \appfile{mermaid/} & When does each milestone open, close, and produce its artifact? \\
14 & 5 & Graphviz & \appfile{graphviz/} & What has to fail, and in what combination, to stop the programme? \\
15 & 5 & PlantUML & \appfile{plantuml/} & What do three parties do at once to close one milestone? \\
16 & 6 & D2 & \appfile{d2/} & Which published quantities can a funder check before committing? \\
17 & 6 & Graphviz & \appfile{graphviz/} & Where does each published number land in a filing? \\
18 & 7 & Diagrams & \appfile{diagrams-python/} & Who is employed, who is contracted, who is contributed? \\
19 & 8 & Mermaid & \appfile{mermaid/} & What did the July and August 2026 contacts actually unlock? \\
20 & 10 & Diagrams & \appfile{diagrams-python/} & What survives if the programme stops, and who holds it? \\
\end{tabularx}
\end{apptable}
\tabcap{The twenty figures, their section, their platform, the directory that
carries each specification, and the one question each answers. Five mermaid,
three plantuml, five d2, three diagrams, four graphviz, by purpose not quota.}

\begin{apptable}
\begin{tabularx}{\textwidth}{>{\raggedright\arraybackslash}p{0.7cm} >{\raggedright\arraybackslash}p{1.1cm} >{\raggedright\arraybackslash}p{4.6cm} >{\raggedright\arraybackslash}X}
\headrow \thead{Tab} & \thead{\S} & \thead{Contents} & \thead{The question it answers}\\
\midrule
1 & 0 & Four readers, four entry points & Which section should this reader open first? \\
2 & 0 & Figure index, twenty rows & Where is each figure and what does it answer? \\
3 & 0 & Table index, nineteen rows & Where is each table and what does it answer? \\
4 & 1 & Three clauses against four eligibility tests & Which clause survives the filter, and why? \\
5 & 1 & Three overhead regimes on one direct base & What is the premium for routing through an institution? \\
6 & 2 & The owned register, thirteen rows with DOIs & What exists today and where is it deposited? \\
7 & 2 & The absent register, seven rows with signatories & What is missing and who alone can supply it? \\
8 & 3 & The four-layer budget under two prices & What does the SBIR route reach, and what does it not? \\
9 & 3 & Phase I line items, eight rows to \$306,000 & What does nine months of Phase I actually buy? \\
10 & 3 & Phase II line items, ten rows to \$1,300,000 & What does twenty-four months of Phase II buy? \\
11 & 4 & The \S54.2 and \S54.4 disclosure triggers & What must be disclosed, at what threshold, on which form? \\
12 & 4 & The capitalization table across five stages & Who owns what, and does the SBIR test still pass? \\
13 & 5 & The twelve milestones, costed and dated & What is the milestone, when, and what does it cost? \\
14 & 5 & Artifact, stop condition, funder mechanism & What proves it, what halts it, and who pays? \\
15 & 6 & Six checkable quantities with limitations & What can a reviewer verify against a published source? \\
16 & 6 & Cost and schedule against industry benchmarks & What did the prior work cost, and what does it cost elsewhere? \\
17 & 7 & The 2.6 FTE by employment status & Who is on the payroll and who is not? \\
18 & 9 & Risks, likelihood, and the mitigation held & What could go wrong and what is already in place? \\
19 & 11 & Abbreviations, two column pairs & What does this term mean in this paper? \\
\end{tabularx}
\end{apptable}
\tabcap{The nineteen tables, their section, their contents, and the one question
each answers. Every table is set to the full body measure with ragged-right
fixed columns, so no cell shows a stretched interword gap.}
